\section{Credit modelling}\label{sec:credit}

We treat each apartment as a one-name obligor whose default time
$\tau_i$ depends on a common factor and an idiosyncratic component.
Three canonical specifications are implemented in
\texttt{src/tranche\_pricing/credit/} and run on a shared Monte Carlo
seed.

\subsection{Gaussian one-factor copula}\label{ssec:gauss-copula}

Following \autocite{Li2000, Vasicek2002}, each obligor's latent asset value
$X_i$ is
\begin{equation}
X_i = \sqrt{\rho}\, M + \sqrt{1 - \rho}\, Z_i,
\qquad M, Z_i \overset{\text{i.i.d.}}{\sim} \mathcal{N}(0, 1),
\label{eq:gauss-copula}
\end{equation}
so that $\text{Corr}(X_i, X_j) = \rho$ for $i \neq j$ and each marginal is
standard normal. Default times are derived under constant hazard
$h = -\log(1 - p) / T$ from the uniform marginal $U_i = \Phi(X_i)$:
$\tau_i = -\log(1 - U_i) / h$, capped at $\infty$ when $U_i \ge p$. The
conditional default probability given the factor is
\begin{equation}
\Prob(\tau_i \le t \given M = m) = \Phi\!\left(
\frac{\Phi^{-1}(1 - e^{-h t}) - \sqrt{\rho}\, m}{\sqrt{1 - \rho}}
\right),
\label{eq:cond-pd}
\end{equation}
which underlies the Vasicek (2002) large-portfolio limit we use as a
control-variate for the senior tranche.

\subsection{Student-t copula}\label{ssec:t-copula}

The same factor structure is mixed with an inverse-chi-square scale to
deliver $t_\nu$ marginals while preserving the linear correlation $\rho$:
\begin{equation}
X_i = \bigl(\sqrt{\rho}\, M + \sqrt{1 - \rho}\, Z_i\bigr) \cdot
\sqrt{\nu / W}, \quad W \sim \chi^2_{\nu}.
\label{eq:t-copula}
\end{equation}
The Gaussian copula is recovered as $\nu \to \infty$. For finite $\nu$ the
Student-t copula features non-zero lower tail dependence,
\begin{equation}
\lambda_L(\rho, \nu) = 2\, t_{\nu + 1}\!\left(-\sqrt{(\nu + 1)\frac{1 - \rho}{1 + \rho}}\right),
\label{eq:tail-dep}
\end{equation}
which \cref{fig:tail-dep} plots over a grid of $(\rho, \nu)$. The Gaussian
copula has $\lambda_L \equiv 0$ for any $\rho < 1$; for $\nu = 5$ and
$\rho = \rhoBaseline{}$ the t-copula carries
$\lambda_L \approx 0.08$, a difference that propagates directly into the
senior tranche's loss tail.

\subsection{Cox doubly-stochastic intensity}\label{ssec:cox}

The asset-value copulas above are static: default correlation comes
from a single Gaussian draw of $M$ at $t = 0$. A doubly-stochastic
intensity model lets the hazard rate vary over time through a macro
factor. We adopt
\begin{equation}
\lambda_i(t) = \alpha + \beta\, Y_t, \qquad
\mathrm{d}Y_t = \kappa (\theta - Y_t) \dt + \xi \sqrt{Y_t}\, \dW^Y_t,
\label{eq:cox}
\end{equation}
with $Y_0 = \theta$ and the CIR positivity preserved by Euler full
truncation. Default times are $\tau_i = \inf\{t : \int_0^t \lambda_i(s)\,
\mathrm{d}s > E_i\}$, $E_i \sim \text{Exp}(1)$. Conditional on the path of
$Y$, defaults are independent.

\subsection{Vasicek (2002) large-portfolio limit}\label{ssec:vasicek-lpf}

For sufficiently large portfolios under the Gaussian one-factor copula
the cumulative loss converges (Vasicek 2002) to a deterministic function
of the common factor:
\begin{equation}
\Prob(L \le \ell) =
\Phi\!\left(\frac{\sqrt{1 - \rho}\,\Phi^{-1}(\ell) - \Phi^{-1}(p)}{\sqrt{\rho}}\right).
\label{eq:vasicek-lpf}
\end{equation}
We use it (a) as a control variate for the senior tranche expected loss
in the variance-reduction layer and (b) as a sanity-check benchmark for
the MC estimator at the largest sample sizes reported in
\cref{fig:mc-convergence}.

\subsection{Student-t tail dependence derivation}\label{ssec:t-dep-derivation}

Following \autocite{EmbrechtsMcNeilStraumann2002}, the lower tail
dependence of a bivariate elliptical distribution with correlation $\rho$
and degrees of freedom $\nu$ is, by definition,
\begin{equation*}
\lambda_L = \lim_{q \to 0^+}\Prob\bigl(U_2 \le q \mid U_1 \le q\bigr).
\end{equation*}
For the Student-t copula, the conditional Student-t representation
yields the closed form quoted in \cref{eq:tail-dep}: the radial structure
of the mixing scale ensures non-zero $\lambda_L$ for all $\nu < \infty$.
The Gaussian copula is recovered in the limit $\nu \to \infty$ and has
$\lambda_L \equiv 0$, the empirical critique of
\autocite{DonnellyEmbrechts2010}.

\subsection{Cox doubly-stochastic survival probability}\label{ssec:cox-survival}

Conditional on the path of the macro factor $Y$, the survival probability
to time $T$ of obligor $i$ is
\begin{equation*}
S_i(T \mid Y) = \exp\!\Big(-\int_0^T \lambda_i(s)\,\mathrm{d}s\Big).
\end{equation*}
For the CIR-affine parameterisation $\lambda_i(t) = \alpha + \beta Y_t$
with $Y_t$ a CIR process, the unconditional survival probability is
itself affine:
\begin{equation*}
S_i(T) = \exp\bigl(-\alpha T - \beta\,A(T)\bigr) \cdot \E\bigl[\exp(-\beta \cdot \text{integrated CIR})\bigr],
\end{equation*}
which has a known Laplace transform closed form (see
\autocite{Schonbucher2003}, Chapter 7). We use it as a sanity test on the
calibration: at the configured $(\kappa, \theta, \xi)$ the closed-form
survival probability at $T = 10$ should agree with the empirical
$1 - $ default-rate in our Monte Carlo to within MC noise.

\subsection{Numerical comparison}

\Cref{fig:loss-dist} overlays the simulated cumulative-loss histograms
under each model at the Paris-intermediate calibration. All three deliver
the same mean loss (matched by construction to the marginal default
probability) but differ markedly in the right tail. The implications for
tranche fair value are quantified in \cref{sec:results}.

\begin{figure}[ht]
  \centering
  \begin{minipage}{0.48\textwidth}
    \includegraphics[width=\linewidth]{fig_tail_dependence.pdf}
    \caption{Lower tail dependence $\lambda_L$ of the Student-t copula as a
    function of $\rho$ for four values of $\nu$. Gaussian copula
    ($\lambda_L \equiv 0$) plotted for reference.}
    \label{fig:tail-dep}
  \end{minipage}\hfill
  \begin{minipage}{0.48\textwidth}
    \includegraphics[width=\linewidth]{fig_loss_distributions.pdf}
    \caption{Cumulative portfolio-loss density under the three credit
    models. Vertical dashed lines mark each model's empirical $\VaR_{95}$.}
    \label{fig:loss-dist}
  \end{minipage}
\end{figure}
